\documentclass[reqno, answers]{exam}

\usepackage{amssymb, amsmath, amsthm, stmaryrd}
\usepackage{fullpage, parskip}
\usepackage{enumitem}
\usepackage{tikz}
\usepackage{ulem}
\usepackage{hyperref}

\title{Introduction to Computational Math Spring 2026 \\ Homework 05}
\date{Due: Friday, Feb 27, 2026, 11:59 PM}
\author{}

\begin{document}
\maketitle
\thispagestyle{empty}

Textbook = Driscoll and Braun (\url{https://fncbook.com/})

There are more problems in this homework than the previous one, but they are shorter and more conceptual.


\begin{enumerate}

%% Exercise 3.2.1
\item Textbook Exercise 3.2.1: Work out the least-squares solution when
\[
A = \begin{bmatrix} 2 & -1 \\ 0 & 1 \\ -2 & 2 \end{bmatrix}, \quad b = \begin{bmatrix} 1 \\ -5 \\ 6 \end{bmatrix}.
\]

\begin{solution}
$x = \begin{bmatrix} -2 \\ -1 \end{bmatrix}$
\end{solution}

%% Exercise 3.2.2
\item Textbook Exercise 3.2.2: Use equation (3.2.4) to find the pseudoinverse $A^+$ of the matrix $A = \begin{bmatrix} 1 & -2 & 3 \end{bmatrix}^T$.

\begin{solution}
$A^+ = \begin{bmatrix} \frac{1}{14} & -\frac{1}{7} & \frac{3}{14} \end{bmatrix}$
\end{solution}

%% Exercise 3.2.3
\item Textbook Exercise 3.2.3: Prove the first statement of Theorem 3.2.2: $A^T A$ is symmetric for any $m \times n$ matrix $A$ with $m \geq n$.

\begin{solution}
$(A^T A)^T = A^T (A^T)^T = A^T A$
\end{solution}

%% Exercise 3.2.5a
\item Textbook Exercise 3.2.5a: Show that for any $m \times n$ matrix $A$ with $m > n$ for which $A^T A$ is nonsingular, $A^+ A$ is the $n \times n$ identity.

\begin{solution}
$A^+ A = (A^T A)^{-1} A^T A = (A^T A)^{-1} (A^T A) = I_n$
\end{solution}

%% Exercise 3.4.1
\item Textbook Exercise 3.4.1: Find a Householder reflector $P$ such that
\[
P \begin{bmatrix} 2 \\ 9 \\ -6 \end{bmatrix} = \begin{bmatrix} 11 \\ 0 \\ 0 \end{bmatrix}.
\]

\begin{solution}
$P = \frac{1}{11} \begin{bmatrix} 2 & 9 & -6 \\ 9 & 2 & 6 \\ -6 & 6 & 7 \end{bmatrix}$
\end{solution}

%% Exercise 3.4.2
\item Textbook Exercise 3.4.2: Prove the unfinished items in Theorem 3.4.1, namely, that a Householder reflector $P$ is symmetric and orthogonal.

To show that a square matrix $Q$ is orthogonal, it is enough to show that $Q^T Q = I$.

\begin{solution}
$P^T = (I - 2\frac{vv^T}{v^T v})^T = I - 2\frac{vv^T}{v^T v} = P$ (symmetric). 

$P^2 = (I - 2\frac{vv^T}{v^T v})^2 = I - 4\frac{vv^T}{v^T v} + 4\frac{vv^T}{v^T v} = I$ (orthogonal).
\end{solution}

%% Exercise 3.4.3
\item Textbook Exercise 3.4.3: Let $P$ be a Householder reflector as in (3.4.1).
\begin{enumerate}
\item[(a)] Find a vector $u$ such that $Pu = -u$.
\item[(b)] What algebraic condition is necessary and sufficient for a vector $x$ to satisfy $Px = x$? In $n$ dimensions, how many such linearly independent vectors are there?
\end{enumerate}

\begin{solution}
\textbf{(a)} $u = v$ (or any scalar multiple of $v$), since $Pv = v - 2v = -v$.

\textbf{(b)} Condition: $v^T x = 0$ (i.e., $x$ orthogonal to $v$). In $n$ dimensions, there are $n-1$ linearly independent such vectors.
\end{solution}

%% Exercise 3.4.7a
\item Textbook Exercise 3.4.7a: Given a 2-vector $[\alpha, \beta]^T$, a Givens rotation defines an angle $\theta$ such that
\[
\begin{bmatrix} \cos\theta & \sin\theta \\ -\sin\theta & \cos\theta \end{bmatrix} \begin{bmatrix} \alpha \\ \beta \end{bmatrix} = \begin{bmatrix} \sqrt{\alpha^2 + \beta^2} \\ 0 \end{bmatrix}.
\]
Given $\alpha$ and $\beta$, show how to compute $\cos\theta$ and $\sin\theta$ without evaluating any trig functions.

This question is asking you to find $\theta$ in terms of $\alpha$ and $\beta$. Your answer will involve a trigonometric function.

\begin{solution}
Let $r = \sqrt{\alpha^2 + \beta^2}$. Then:
\[
\cos\theta = \frac{\alpha}{r}, \quad \sin\theta = \frac{\beta}{r}
\]
\end{solution}

\end{enumerate}

No Jupyter notebook this week.



\section*{Quiz Information}

The quiz will be held on {\bf Tuesday, March 03, 2026} during the discussion section.

Quiz questions will be modifications of the homework problems and material discussed in class.

\textbf{Topics covered:} 

Chapters 3.2 and 3.4. 

You will need memorize 
\begin{itemize}
    \item The normal equations for least squares problems,
    \item The formula for the Householder reflector.
\end{itemize}

You should expect both computational and conceptual questions on the quiz.

There won't be any questions about using the normal equations in terms of $QR$ factorization, but you should know how to use the normal equations in terms of the $A^T A$ matrix.



\end{document}